\chapter{Additional Probability and Asymptotic Tools}
\label{app:probability}

This appendix provides a more careful probability-space description, a few
distributional examples, simulation code for the WLLN and CLT, and the delta
method.  None of it is required on a first reading of
Chapter~\ref{ch:probability}.

\section{Probability spaces, events, and random variables}

A probability space is a triple $(\mathcal S,\mathcal F,P)$:
\begin{itemize}
  \item $\mathcal S$ is the \emph{sample space} of elementary outcomes;
  \item $\mathcal F$ is a sigma-algebra of events, meaning it contains
        $\mathcal S$ and is closed under complements and countable unions;
  \item $P:\mathcal F\to[0,1]$ satisfies $P(\mathcal S)=1$ and countable
        additivity over pairwise disjoint events.
\end{itemize}
A real random variable is a measurable map
$X:\mathcal S\to\mathbb R$.  Its distribution is the probability measure
induced on the real line:
\[
  P_X(B)=P\{s\in\mathcal S:X(s)\in B\}.
\]
Thus, the sample space contains underlying outcomes, whereas the support of
$X$ contains possible numerical values.  Treating those two sets as the same
is a convenient shortcut in elementary examples, but it is not the general
definition.

For example, an outcome can record tomorrow's temperature, rainfall, and
wind.  The random variable $X$ may retain only the indicator that it rains.
Then the sample space is rich, while the support of $X$ is merely $\{0,1\}$.

\section{CDFs, densities, and continuous distributions}

Every real random variable has a right-continuous c.d.f.
$F_X(x)=P(X\leq x)$.  A variable is \emph{continuous} (atomless) if
$P(X=x)=0$ for every $x$, equivalently if its c.d.f. is continuous.  It is
\emph{absolutely continuous} if there exists an integrable density $f_X$
such that
\[
  F_X(x)=\int_{-\infty}^{x}f_X(s)\,ds.
\]
Every absolutely continuous distribution is continuous, but the converse is
false.  The Cantor distribution is the standard counterexample: it has a
continuous c.d.f. but no density with respect to ordinary length (Lebesgue
measure).

The geometric ``probability equals length'' picture is therefore not a
general definition.  If $X$ is uniform on $[0,1]$, then
$P(X\in A)$ equals the length of $A$ for suitable sets $A$.  Under a
nonuniform density, probability is instead weighted length,
$P(X\in A)=\int_Af_X(x)\,dx$; for a discrete distribution, it is a sum of
point masses.

The code below compares three normal densities.  It uses base R so that the
live page loads quickly.

\begin{rcode}
x <- seq(-7, 7, length.out = 400)
plot(x, dnorm(x, 0, 1), type = "l", lwd = 2,
     xlab = "x", ylab = "density", ylim = c(0, 0.42))
lines(x, dnorm(x, 2, 1), lty = 2, lwd = 2, col = "steelblue")
lines(x, dnorm(x, 0, 2), lty = 3, lwd = 2, col = "firebrick")
legend("topright", c("N(0,1)", "N(2,1)", "N(0,4)"),
       lty = 1:3, lwd = 2,
       col = c("black", "steelblue", "firebrick"), bty = "n")
\end{rcode}

\section{Conditional distributions with continuous variables}

When $P(X=x)=0$, elementary event conditioning cannot define
$P(Y\leq y\mid X=x)$ by a ratio of probabilities.  A regular conditional
distribution supplies a mathematically valid version.  When a joint density
exists and $f_X(x)>0$, it has the familiar formula
\[
  f_{Y\mid X}(y\mid x)=\frac{f_{X,Y}(x,y)}{f_X(x)}.
\]
One sometimes motivates this expression using limits of shrinking intervals
around $x$.  Such a limit is an intuition, not a universal definition: its
validity requires regularity conditions and it identifies the conditional
law only for almost every $x$.

Bayes' rule for densities is
\[
  f_{X\mid Y}(x\mid y)
  =\frac{f_{Y\mid X}(y\mid x)f_X(x)}{f_Y(y)},
  \qquad f_Y(y)>0.
\]
Marginalization gives
\[
  f_X(x)=\int f_{X,Y}(x,y)\,dy,
  \qquad
  F_X(x)=\int F_{X\mid Y}(x\mid y)\,dF_Y(y).
\]

\section{The bivariate normal example}

A random vector $(X,Y)'$ is jointly normal if every linear combination
$aX+bY$ is normally distributed.  For a jointly normal pair, zero covariance
is equivalent to independence.  This equivalence is special: it is false for
general distributions.

The following optional demonstration plots bivariate normal densities with
positive, negative, and zero correlation.  The live version uses
\texttt{mvtnorm} and \texttt{plotly}; the two-dimensional contour option is
usually quicker in class.

\begin{rcode}
library(mvtnorm)
x <- seq(-4, 4, length.out = 80)
y <- seq(-4, 4, length.out = 80)
mu <- c(0, 0)
Sigma <- matrix(c(1, 0.5, 0.5, 1), 2, 2)
z <- outer(x, y, Vectorize(function(a, b) {
  dmvnorm(c(a, b), mean = mu, sigma = Sigma)
}))
contour(x, y, z, xlab = "X", ylab = "Y")
\end{rcode}

\section{A compact WLLN and CLT simulation}

The original classroom simulation repeated similar commands for ten sample
sizes.  The vectorized version below does the same job more transparently.
For Bernoulli$(0.5)$ observations, the WLLN predicts concentration of
$\bar X_n$ around $0.5$, and the CLT predicts
\[
  \frac{\sqrt n(\bar X_n-0.5)}{0.5}\dto N(0,1),
\]
because the Bernoulli standard deviation is $0.5$.

\begin{rcode}
set.seed(12)
B <- 1000
n_grid <- c(10, 25, 50, 100, 250)

sim_one_n <- function(n) {
  means <- replicate(B, mean(rbinom(n, 1, 0.5)))
  data.frame(
    n = n,
    mean = means,
    standardized = sqrt(n) * (means - 0.5) / 0.5
  )
}
sim <- do.call(rbind, lapply(n_grid, sim_one_n))

# WLLN: estimated probability of falling within 0.1 of the mean
aggregate(abs(mean - 0.5) < 0.1 ~ n, sim, mean)

# CLT: compare one standardized histogram with N(0,1)
z <- subset(sim, n == max(n_grid))$standardized
hist(z, probability = TRUE, breaks = 25,
     main = paste("n =", max(n_grid)), xlab = "standardized mean")
curve(dnorm(x), add = TRUE, lwd = 2, col = "firebrick")
\end{rcode}

\section{The delta method}

The delta method transfers an asymptotic distribution through a smooth
transformation.  It is a first-order Taylor expansion with probabilistic
remainder control.

\begin{theorem}[Multivariate delta method]
Suppose $X_n\in\mathbb R^k$ satisfies
\[
  \sqrt n(X_n-\theta)\dto N(0,\Sigma),
\]
and $g:\mathbb R^k\to\mathbb R^\ell$ is continuously differentiable near
$\theta$.  Let $G(\theta)$ be the $\ell\times k$ Jacobian whose $(r,j)$ entry
is $\partial g_r(\theta)/\partial\theta_j$.  Then
\[
  \sqrt n\{g(X_n)-g(\theta)\}
  \dto N\!\left(0,G(\theta)\Sigma G(\theta)'\right).
\]
\end{theorem}

For a scalar-valued $g$, $G(\theta)$ is a row vector and the limiting variance
is $G(\theta)\Sigma G(\theta)'$.  A statement such as ``converges to
$N(0,\cdot)$'' must include the $N$; a covariance matrix alone is not a
distribution.  If the first derivative is zero, this first-order result is
degenerate and a higher-order delta method may be needed.

\begin{example}[Log transformation]
If $\sqrt n(\widehat\theta-\theta)\dto N(0,V)$ with $\theta>0$, then
$g(\theta)=\log\theta$ has derivative $1/\theta$, so
\[
  \sqrt n\{\log\widehat\theta-\log\theta\}
  \dto N\!\left(0,\frac{V}{\theta^2}\right).
\]
\end{example}
